% F17_Question.tex
% Title: The F17 Question: Could a 17th Operator Improve Any SuperBEST Entry?
% Author: Arturo R. Almaguer
% Date: April 2026
% Session: X20
% Status: THEORETICAL ESSAY — open question and analysis
% Companion to: SuperBEST_v5_Structural_Audit.tex, F16_Characterization.tex

\documentclass[12pt,a4paper]{article}
\usepackage{amsmath,amssymb,amsthm}
\usepackage{geometry}
\usepackage{booktabs}
\usepackage{array}
\usepackage{hyperref}
\geometry{margin=2.5cm}

%% ── Theorem environments ──────────────────────────────────────────────────────
\newtheorem{theorem}{Theorem}[section]
\newtheorem{lemma}[theorem]{Lemma}
\newtheorem{corollary}[theorem]{Corollary}
\newtheorem{proposition}[theorem]{Proposition}
\theoremstyle{definition}
\newtheorem{definition}[theorem]{Definition}
\newtheorem{remark}[theorem]{Remark}
\theoremstyle{remark}
\newtheorem{openq}[theorem]{Open Question}

%% ── Operator macros ───────────────────────────────────────────────────────────
\newcommand{\EML}{\operatorname{EML}}
\newcommand{\EXL}{\operatorname{EXL}}
\newcommand{\EAL}{\operatorname{EAL}}
\newcommand{\DEML}{\operatorname{DEML}}
\newcommand{\ELAd}{\operatorname{ELAd}}
\newcommand{\ELSb}{\operatorname{ELSb}}
\newcommand{\ELMl}{\operatorname{ELMl}}
\newcommand{\LEAd}{\operatorname{LEAd}}
\newcommand{\LEdiv}{\operatorname{LEdiv}}

%% ── Function macros ───────────────────────────────────────────────────────────
\newcommand{\recip}{\operatorname{recip}}
\newcommand{\divop}{\operatorname{div}}
\newcommand{\mulop}{\operatorname{mul}}
\newcommand{\subop}{\operatorname{sub}}
\newcommand{\addop}{\operatorname{add}}
\newcommand{\negop}{\operatorname{neg}}
\newcommand{\powop}{\operatorname{pow}}
\newcommand{\sqrtop}{\operatorname{sqrt}}
\newcommand{\SQR}{\operatorname{SQR}}
\newcommand{\SMUL}{\operatorname{SMUL}}
\newcommand{\SIN}{\operatorname{SIN}}
\newcommand{\SB}{\mathrm{SB}}
\newcommand{\F}{\mathcal{F}_{16}}
\newcommand{\Fs}{\mathcal{F}_{17}}
\newcommand{\R}{\mathbb{R}}
\newcommand{\Q}{\mathbb{Q}}
\newcommand{\ELC}{\mathrm{ELC}(\mathbb{R})}
\newcommand{\ELfin}{\mathrm{EL}_{\mathrm{fin}}}

%% =============================================================================
\title{The F17 Question:\\[4pt]
       Could a 17th Operator Improve Any SuperBEST Entry?}
\author{Arturo R.\ Almaguer\\
  \small monogate research\quad\texttt{almaguer1986@gmail.com}}
\date{April 2026 (Session X20)}
%% =============================================================================

\begin{document}
\maketitle

%% ── Abstract ──────────────────────────────────────────────────────────────────
\begin{abstract}
The SuperBEST v5 table records the minimum $\mathcal{F}_{16}$-node counts
for ten standard arithmetic operations, totalling exactly 18 nodes.
All ten entries are proved optimal: no improvement is possible within $\F$.
We ask whether there is a \emph{natural} 17th operator that, when added to
$\F$, would reduce any entry.

The answer is yes, in three ways.
First, \textbf{SQR} (square root as a 1-input primitive) reduces
$\SB(\sqrtop)$ from $2$ to $1$, by the same structural argument that
reduced $\SB(\recip)$ from $2$ to $1$ via $\ELSb$ in the R16 sprint.
Second, if $\ELMl(x,y) = e^{x\ln y} = y^x$ (the Class~I exponentiation
operator listed in the 16-operator census) is admitted as an active
operator, then $\SB(\powop) = 1$ --- a collapse from the current
3-node result.
Third, a ternary multiply $\mathrm{TMUL}(x,y,z)=xyz$ reduces
triple-product computations from $3n$ to $1n$ in the catalog,
though it leaves the core table unchanged.

None of these extensions enlarges $\ELC$: they are all already within the
exp-log closure.  The only extension that strictly enlarges $\ELC$ is adding
$\sin$ as a primitive (the Richardson-elementary extension), which changes the
research problem entirely.

An \textbf{internal consistency question} emerges: if $\ELMl$ is
listed among the 16 operators in the census but $\SB(\powop)=3$ in
SuperBEST~v5, then either (a)~$\ELMl$ is not in the $\F$ subset used
by v5, or (b)~the $\powop$ entry has an error.
Resolution is the highest-priority open question identified in this session.
\end{abstract}

\tableofcontents

%% =============================================================================
\section{Setup: What Does the SuperBEST v5 Table Actually Say?}
\label{sec:setup}
%% =============================================================================

\subsection{The Table}

SuperBEST v5 \cite{SBv5} proves that ten core operations require exactly
the following numbers of $\F$-nodes:

\begin{table}[h]
\centering
\caption{SuperBEST v5 (all entries proved optimal)}
\label{tab:v5}
\begin{tabular}{@{}lccl@{}}
\toprule
\textbf{Operation} & \textbf{Cost} & \textbf{Naive} & \textbf{Lower Bound Method} \\
\midrule
$\exp(x)$        & 1 &  1 & Leaf scan \\
$\ln(x)$         & 1 &  3 & Leaf scan \\
$\recip(x)$      & 1 &  5 & Derivative ($-1/x^2$) \\
$\divop(x,y)$    & 2 & 15 & Derivative obstruction \\
$\negop(x)$      & 2 &  9 & Slope barrier ($-1$) \\
$\mulop(x,y)$    & 2 & 13 & Derivative obstruction (T32) \\
$\subop(x,y)$    & 2 &  5 & Derivative obstruction (T33) \\
$\addop(x,y)$    & 2 & 11 & Slope barrier ($+1$, ADD-T1) \\
$\sqrtop(x)$     & 2 &  8 & Derivative ($1/(2\sqrt{x})$) \\
$\powop(x,n)$    & 3 & 15 & 2-node enumeration \\
\midrule
\textbf{Total}   & \textbf{18} & \textbf{73} & \\
\bottomrule
\end{tabular}
\end{table}

\subsection{All Entries Are Tight}

\begin{theorem}[SuperBEST v5 Optimality, \cite{SBv5}]
Every entry in Table~\ref{tab:v5} is exact: $\SB(f) = k$ iff the
construction achieves $k$ nodes and the lower-bound argument rules out $k-1$.
No entry can be improved within $\F$.
\end{theorem}

Consequently, there are no ``gaps'': the question is not ``which entry has
the largest gap?'' but rather ``which entry is most \emph{vulnerable} to
a natural operator extension?''

The answer by inspection is $\powop = 3$, the only 3-node entry.  Every
other entry costs $\le 2$ nodes.  Any F17 that reduces $\powop$ saves
2 nodes and brings the table total to 16.  Any F17 that reduces $\sqrtop$
saves 1 node and brings the total to 17.

%% =============================================================================
\section{The Multiplication Barrier and Signed Inputs}
\label{sec:mul-barrier}
%% =============================================================================

The mul entry $\SB(\mulop) = 2$ is proved with the domain restriction $x > 0$.
The 2-node construction $\ELAd(\EXL(0,x),\,y) = e^{\ln x}\cdot y = xy$
requires $\ln x$, which is undefined for $x \le 0$.

For general signed multiplication, the F16 route requires a negation detour:
\begin{align*}
  xy &= -\bigl((-x)\cdot y\bigr) \quad\text{for } x < 0 \\
  &= -\bigl(\ELAd(\EXL(0,-x),\,y)\bigr) \\
  &= \negop\!\bigl(\ELAd(\EXL(0,\negop(x)),\,y)\bigr),
\end{align*}
costing $2 + 2 = 4$ additional nodes for the negation steps.

\begin{proposition}[F17a: Signed Multiply Primitive]
\label{prop:smul}
If a primitive $\SMUL(x,y) = xy$ for all $x,y\in\R$ is added to $\F$,
then $\SB_{\Fs}(\mulop_{\mathrm{gen}}) = 1$.
However, $\SMUL$ is not a composition of $e^{\pm x}$ and $\ln y$ via any
arithmetic operation.  It would constitute a genuinely new computational
primitive, architecturally alien to $\F$.
\end{proposition}

\begin{proof}
The first claim is immediate: $\SMUL(x,y) = xy$ is the desired function,
computed in 1 node by definition.

The second claim follows by inspection of the $\F$ construction:
every $\F$ operator has the form $h(e^{\varepsilon u}, \ln v)$ or
$h(e^{\varepsilon u}, v)$ for $h\in\{+,-,\times,\div\}$.
None of these equals $u\cdot v$ for all $u,v\in\R$: the closest,
$\ELAd(u,v) = e^u\cdot v$, differs by the exponential factor $e^u \ne u$.
\end{proof}

\textbf{Conclusion:} F17a = SMUL is possible but breaks the structural
homogeneity of $\F$ (every $\F$ operator wraps arithmetic in exp or ln).
It is not a natural extension.

%% =============================================================================
\section{The Absolute Value Barrier}
\label{sec:abs}
%% =============================================================================

\begin{theorem}[\cite{F16Char}]
$|x|\notin\ELfin$.
\end{theorem}

\begin{proof}[Proof sketch]
Every $\F$ tree computes a real-analytic function on its natural domain.
The function $|x|$ is not differentiable at $x=0$.  Hence no finite $\F$
tree computes $|x|$.
\end{proof}

\noindent Adding $|x|$ as a primitive would extend $\ELfin$ beyond $\ELC$ ---
the set of representable functions would include non-analytic functions.
This is a fundamentally different extension from those discussed elsewhere.

%% =============================================================================
\section{F17b: Square Root as a Primitive}
\label{sec:sqrt-primitive}
%% =============================================================================

\begin{definition}[F17b]
$\mathcal{F}_{17b} = \F \cup \{\SQR\}$, where $\SQR(x) = \sqrt{x}$ for $x>0$.
\end{definition}

\begin{theorem}[SQR reduces $\sqrtop$ to 1 node]
\label{thm:sqr}
$\SB_{\mathcal{F}_{17b}}(\sqrtop) = 1$.
All other SuperBEST v5 entries are unchanged.
The new table total is $\mathbf{17}$ nodes.
\end{theorem}

\begin{proof}
\textbf{Upper bound:} $T = \SQR(x)$ computes $\sqrt{x}$ in 1 node.

\textbf{Lower bound (unchanged):} The lower bound $\SB(\sqrtop)\ge 2$
in~\cite{SBv5} applies specifically within $\F$.  Once $\SQR\in\Fs$,
the 1-node construction achieves $\SB_{\Fs}(\sqrtop) = 1$.

\textbf{Other entries unchanged:} No other entry uses $\SQR$ as an
intermediate.  The lower bounds for the other 9 entries are proved by
derivative or slope arguments that are independent of whether $\SQR$ is
in the operator library.
\end{proof}

\subsection{Structural Motivation}

This extension has a clear historical precedent within the monogate project.
In the R16 sprint \cite{R16}, the operator $\ELSb(x,y) = e^x/y$ was
recognized as encoding $\recip(x) = 1/x$ in 1 node (via $\ELSb(0,x)$),
reducing $\SB(\recip)$ from $2$ (in $\mathcal{F}_6$) to $1$ (in $\mathcal{F}_{16}$).
The argument was: $\recip$ is already in $\ELC$; $\ELSb$ is the exp-ln
composition that computes it; admitting $\ELSb$ as a primitive reduces the cost.

The same argument applies to $\SQR$:

\begin{itemize}
  \item $\sqrt{x}\in\ELC$ (it equals $e^{(1/2)\ln x}$).
  \item The 2-node $\F$ route $\EML\!\bigl(\tfrac{1}{2}\EXL(0,x),\,1\bigr)$
    computes $\sqrt{x}$.
  \item If $\SQR$ is added as a primitive, the cost drops to 1.
  \item $\SQR$ is analytic for $x>0$, preserving the analyticity of $\Fs$.
  \item $\ELfin = \ELC$ remains valid: $\SQR(x) = e^{(1/2)\ln x}\in\ELC$,
    so $\Fs$ still computes only $\ELC$ functions.
\end{itemize}

\begin{remark}[SQR vs.\ ELSb analogy]
$\SQR(x) = e^{(1/2)\ln x}$ is to $\ELSb(0,x) = e^{0-\ln x} = 1/x$ as
half-power is to negative-one-power.  Both are specializations of
$e^{c\ln x} = x^c$ for rational $c$.  The pattern suggests a more general
extension: admitting $x^c$ for any fixed rational $c$ as a 1-node primitive.
This is the ``rational power'' extension, which would make all entries of the
form $x^{p/q}$ cost 1 node.
\end{remark}

%% =============================================================================
\section{F17d: The ELMl Question and the Pow Inconsistency}
\label{sec:elml}
%% =============================================================================

\subsection{ELMl in the 16-Operator Census}

The \emph{Sixteen Operator Census} \cite{Census} lists $\ELMl(x,y) = e^{x\ln y} = y^x$
as one of the 16 Class~I operators (with $h = \hat{}\,$ representing
exponentiation).

\begin{definition}[$\ELMl$]
$\ELMl(x,y) = \exp\!\bigl(x\cdot\ln y\bigr) = y^x$, defined for $y > 0$.
\end{definition}

Applied to the leaves $n$ and $x$ (with $x > 0$):
\[
  \ELMl(n,\,x) = x^n.
\]
This is a 1-node computation of $\powop(x,n)$.

\subsection{The Inconsistency}

SuperBEST v5 \cite{SBv5} states $\SB(\powop) = 3$, with construction:
\[
  v_1 = \EXL(0,n) = \ln n, \quad
  v_2 = \EXL(v_1,\,x) = n\ln x, \quad
  v_3 = \EML(v_2,\,1) = e^{n\ln x} = x^n.
\]

If $\ELMl\in\F$, then $\ELMl(n,x) = x^n$ in 1 node, contradicting $\SB(\powop) = 3$.

\begin{proposition}[Resolution]
\label{prop:elml-resolution}
The most likely resolution is that SuperBEST v5 uses a restricted
8-operator subset of the 16-operator census:
\[
  \F_{\mathrm{v5}} = \{\EML, \DEML, \EXL, \EAL, \LEAd, \ELAd, \ELSb, \LEdiv\},
\]
which does \emph{not} include $\ELMl$.
Under this restricted definition, $\SB(\powop) = 3$ is correct.
If $\ELMl$ were admitted, $\SB(\powop) = 1$.
\end{proposition}

\begin{remark}[Why exclude ELMl?]
$\ELMl(x,y) = y^x$ requires both inputs to participate in a product inside
the exponential: $e^{x\ln y}$.  This is qualitatively different from the
other operators, which have one exp-or-ln outer function applied to an
arithmetic combination.  The combining operation for $\ELMl$ is
$x\mapsto x\cdot\ln y$ (scaling by a log), rather than $h(e^{\sigma x}, \ln y)$
for a fixed $h$.  In other words, $\ELMl$ is ``two logs deep'' in some sense ---
it is the composition of multiplication with two log-type operations.
The authors of v5 may have excluded $\ELMl$ precisely because admitting it
would trivialize $\powop$ to 1 node and collapse an interesting structural question.
\end{remark}

\begin{openq}
Is $\ELMl$ in $\mathcal{F}_{16}$ as used by SuperBEST v5?
If yes, should the pow entry be revised to $\SB(\powop) = 1$?
If no, what is the precise structural criterion that excludes $\ELMl$
from the active 16-operator set, and how many of the 16 census operators
are actually used in the SuperBEST routing?
\end{openq}

\subsection{F17d Proposal}

\begin{definition}[F17d]
$\mathcal{F}_{17d} = \F \cup \{\ELMl\}$ (adding $\ELMl$ to the v5 active set).
\end{definition}

\begin{theorem}
$\SB_{\mathcal{F}_{17d}}(\powop) = 1$.
The SuperBEST table total under $\mathcal{F}_{17d}$ is $\mathbf{16}$ nodes.
\end{theorem}

\begin{proof}
$\ELMl(n,x) = x^n$ computes $\powop$ in 1 node.
The lower bound $\SB(\powop)\ge 2$ proved in~\cite{SBv5} applies only within
$\F$ (not $\mathcal{F}_{17d}$): specifically, it checks all 2-node trees over
$\F$-operators and finds none computing $x^n$.  Once $\ELMl$ is admitted,
the 1-node tree $\ELMl(n,x)$ is available.
The remaining 9 entries are unchanged (their lower bounds do not depend on
whether $\ELMl$ is in the library).
\end{proof}

%% =============================================================================
\section{F17c: Ternary Multiplication}
\label{sec:ternary}
%% =============================================================================

\begin{definition}[F17c]
$\mathrm{TMUL}(x,y,z) = x\cdot y\cdot z$ (a ternary operator).
\end{definition}

Within $\F$ (binary operators), computing $xyz$ requires two binary
multiplications:
\[
  xyz = \mulop\!\bigl(\mulop(x,y),\,z\bigr).
\]
By SuperBEST v5, $\mulop = 2$ nodes and the outer $\mulop$ uses the
intermediate $\mulop(x,y)$ as input, costing $2+2 = 4$ nodes total
(the intermediate $\ln(\mulop(x,y))$ is another node).

More precisely, the F16 route for $xyz$:
\begin{align*}
  v_1 &= \EXL(0,x) = \ln x, \\
  v_2 &= \ELAd(v_1,\,y) = xy, \\
  v_3 &= \EXL(0,v_2) = \ln(xy) = \ln x + \ln y, \\
  v_4 &= \ELAd(v_3,\,z) = (xy)\cdot z = xyz.
\end{align*}
Total: 4 nodes.

With $\mathrm{TMUL}$: 1 node.  Savings: 3 nodes per triple product.

\begin{remark}[Catalog impact]
Many physical laws are three-factor products:
Lorentz force $F = qvB\sin\theta$ (4 factors, 3 muls in F16 = 6 nodes;
with TMUL = 2 nodes),
work $W = Fd\cos\theta$ (3 factors = 3 nodes in F16; with TMUL = 1 node),
thermal expansion $\Delta L = \alpha L\,\Delta T$ (3 factors).
The catalog savings depend on the frequency of triple products in physics equations.
This is a catalog-level improvement, not a SuperBEST table improvement
(since TMUL leaves the core 10-entry table unchanged).
\end{remark}

%% =============================================================================
\section{F17e: Adding Sin --- The Richardson Extension}
\label{sec:sin}
%% =============================================================================

\begin{definition}[F17e]
$\mathcal{F}_{17e} = \F \cup \{\SIN\}$, where $\SIN(x) = \sin(x)$.
\end{definition}

\begin{theorem}[\cite{F16Char}]
$\sin(x)\notin\ELC$ and hence $\sin(x)\notin\ELfin$.
\end{theorem}

Adding $\SIN$ to $\F$ strictly extends the representable class:
$\ELfin(\mathcal{F}_{17e}) \supsetneq \ELC$.
This is the \emph{Richardson-elementary extension}, studied by
Richardson~\cite{Richardson1968}.

\begin{remark}[Why F17e defeats the purpose]
The value of $\F$ lies in the characterization $\ELfin = \ELC$.
SuperBEST v5 is meaningful because it quantifies complexity within this
clean class.  Adding $\SIN$ replaces one classification problem (``what
is the minimal $\F$-cost of each $\ELC$ function?'') with a harder one
(``what is the minimal $\mathcal{F}_{17e}$-cost of each Richardson-elementary
function?'').  The latter is open and much harder.
F17e is not an improvement to the $\F$ framework; it is a different framework.
\end{remark}

%% =============================================================================
\section{Summary and Recommendations}
\label{sec:summary}
%% =============================================================================

\begin{table}[h]
\centering
\caption{F17 candidates evaluated}
\label{tab:candidates}
\begin{tabular}{@{}p{3cm}p{2.8cm}ccp{3.5cm}@{}}
\toprule
\textbf{Candidate} & \textbf{Operator} & \textbf{Extends ELC?} &
  \textbf{Table Savings} & \textbf{Verdict} \\
\midrule
F17b & $\SQR(x) = \sqrt{x}$ & No & $-1\,\text{n}$ (total: 17) &
  Natural, analytic, clean analogy to ELSb$\to$recip \\[4pt]
F17d & $\ELMl(x,y) = y^x$ & No & $-2\,\text{n}$ (total: 16) &
  Most impactful; contingent on ELMl status in F16 \\[4pt]
F17a & $\SMUL(x,y) = xy$ (signed) & No & $0$ &
  Architecturally alien; not exp-ln \\[4pt]
F17c & $\mathrm{TMUL}(x,y,z) = xyz$ & No & $0$ &
  Catalog improvement; breaks binary-operator model \\[4pt]
F17e & $\SIN(x)$ & \textbf{Yes} & varies &
  Changes research problem; not recommended \\
\bottomrule
\end{tabular}
\end{table}

\subsection{Main Findings}

\begin{enumerate}
  \item \textbf{SuperBEST v5 is optimal within $\F$ as defined.}  All 10 entries
    are proved exact.  No improvement is possible without extending $\F$.

  \item \textbf{The most natural F17 is $\SQR$.}  Adding $\SQR(x) = \sqrt{x}$
    as a primitive reduces $\SB(\sqrtop) = 2 \to 1$, bringing the table total
    to 17 nodes.  The extension is analytic, within $\ELC$, and follows the
    same pattern as the R16 sprint (ELSb $\to$ recip $= 1$~n).

  \item \textbf{The highest-impact F17 is $\ELMl$}, reducing pow from 3 to
    1 and the table total from 18 to 16.  But whether $\ELMl$ is already in
    $\F$ (and therefore the pow entry has an error) or is genuinely excluded
    (in which case F17d is a meaningful extension) is the key open question.

  \item \textbf{No natural F17 extension enlarges $\ELC$.}  All candidates
    within the analytic, exp-ln spirit of $\F$ are already in $\ELC$.
    The $\mathcal{F}_{17}$ extensions change cost bounds, not representable classes.

  \item \textbf{F16 is the natural complete family for $\ELC$.}  The
    characterization $\ELfin = \ELC$ \cite{F16Char} is the correct boundary.
    F17 extensions within $\ELC$ are refinements of the cost theory, not
    extensions of the representable class.
\end{enumerate}

\subsection{Recommended Follow-up}

\begin{openq}[Priority 1]
Is $\ELMl(x,y) = y^x$ in $\mathcal{F}_{16}$ as used by SuperBEST v5?
Run the 2-node search for pow(x,n) against the v5 operator set.
If $\ELMl$ is admitted, pow $= 1$~n and the table total is 16.
\end{openq}

\begin{openq}[Priority 2]
How many equations in the 315+ catalog contain $\sqrt{x}$?
What is the total catalog node savings if F17b (SQR as primitive) is adopted?
\end{openq}

\begin{openq}[Priority 3]
Define formally the ``rational-power extension'' $\mathcal{F}_{17p}$:
add $x^c$ for every fixed $c\in\Q$ as a primitive.
How does this change the SuperBEST table and the catalog?
\end{openq}

%% =============================================================================
\section*{Conclusion}
%% =============================================================================

The SuperBEST v5 table is complete within $\mathcal{F}_{16}$.
Any improvement requires extending the operator family.
The most natural and architecturally coherent extensions are:

\begin{itemize}
  \item \textbf{F17b} (SQR primitive): reduces $\sqrtop$ from $2$ to $1$,
    following the same R16 pattern that reduced recip.
    New table total: 17 nodes.
  \item \textbf{F17d} (ELMl primitive): if ELMl is not already in $\F$,
    adding it reduces pow from $3$ to $1$.
    New table total: 16 nodes.
\end{itemize}

Both lie within $\ELC$: they do not extend the representable class.
The characterization $\ELfin = \ELC$ remains valid after either extension.

\bigskip
\noindent\textit{Almaguer, A.R.\ (2026).
The F17 Question: Could a 17th Operator Improve Any SuperBEST Entry?
monogate research (Session X20).\ \url{https://monogate.org}}

%% =============================================================================
\begin{thebibliography}{9}
%% =============================================================================

\bibitem{SBv5}
Almaguer, A.R.\ (2026).
\textit{SuperBEST v5: Complete Structural Proof of Optimality for All Ten
Core Arithmetic Operations}.
monogate research.
\texttt{python/paper/theorems/SuperBEST\_v5\_Structural\_Audit.tex}.

\bibitem{F16Char}
Almaguer, A.R.\ (2026).
\textit{The F16 Characterization Theorem: Which Functions Are Exactly
Computable by Finite F16 Trees?}
monogate research.
\texttt{python/paper/theorems/F16\_Characterization.tex}.

\bibitem{Census}
Almaguer, A.R.\ (2026).
\textit{The 16-Operator Census: Completeness Classification of the
Exponential-Logarithmic Binary Operator Family}.
monogate research.
\texttt{python/paper/theorems/Sixteen\_Operator\_Census.tex}.

\bibitem{ADDT1}
Almaguer, A.R.\ (2026).
\textit{General Addition in Two Nodes: Closing the Last Gap in SuperBEST v4}.
monogate research.
\texttt{python/paper/theorems/ADD\_T1\_General\_Addition\_2n.tex}.

\bibitem{R16}
Almaguer, A.R.\ (2026).
\textit{R16: recip $= 1$ node in $\mathcal{F}_{16}$}.
monogate research.
\texttt{python/paper/theorems/recip\_One\_Node.tex}.

\bibitem{Richardson1968}
Richardson, D.\ (1968).
Some unsolvable problems involving elementary functions of a real variable.
\textit{Journal of Symbolic Logic} \textbf{33}(4), 514--520.

\end{thebibliography}

\end{document}
